% Paper-ready section describing the experimental configuration in this repository.
% The image archive and labels_v2.txt are intentionally not tracked, so the absolute
% sample counts cannot be recovered from the repository. Replace the five values below
% with the values printed by split_statistics() before including this file in the paper.
\newcommand{\NumTrainingImages}{\textbf{[insert training total]}}
\newcommand{\NumTrainingHealthy}{\textbf{[insert training healthy]}}
\newcommand{\NumTrainingUnhealthy}{\textbf{[insert training unhealthy]}}
\newcommand{\NumValidationImages}{\textbf{[insert validation total]}}
\newcommand{\NumTestImages}{\textbf{[insert test total]}}

\section{Experimental Setup}
\label{sec:experimental-setup}

This work evaluates whether existing adversarial attack and defense methods are
applicable to the Hydra detector-monitoring workflow at Jefferson Lab.  We do
not propose a new attack or defense algorithm.  Instead, we establish a common
benchmark in which the detector data, classifier, data partition, and evaluation
procedure are held fixed while the attack or defense method is varied.

\subsection{Hydra Detector Data}

Hydra is an AI-assisted monitoring framework that uses detector-generated
images, such as occupancy maps and histogram-like detector signatures, to
provide near-real-time information about experimental health.  We formulate
the archived Hydra data as a binary image-classification task.  Images marked
\emph{good} in the archive are referred to as \emph{healthy} and are assigned
label~1, whereas images marked \emph{bad} are referred to as \emph{unhealthy}
and are assigned label~0 throughout this paper.

The images are resized to $608\times256$ pixels (width $\times$ height),
converted to RGB, and scaled to the interval $[0,1]$.  The archive is divided
once into training, validation, and test sets using a 70:15:15 split, a fixed
random seed of 42, and class-stratified sampling.  Thus, the same partition is
used for classifier training, attack optimization, model selection, and final
comparison.  The training set contains \NumTrainingImages{} images:
\NumTrainingHealthy{} healthy images and \NumTrainingUnhealthy{} unhealthy
images.  The validation and test sets contain \NumValidationImages{} and
\NumTestImages{} images, respectively.  Table~\ref{tab:hydra-splits} reports the
class distribution explicitly so that the imbalance between normal and
abnormal detector conditions is visible rather than being described as
``balanced.''

\begin{table}[t]
    \centering
    \caption{Hydra data partition.  Healthy and unhealthy correspond to the
    archive labels \emph{good} and \emph{bad}, respectively.}
    \label{tab:hydra-splits}
    \begin{tabular}{lrrr}
        \hline
        Split & Healthy (good) & Unhealthy (bad) & Total \\
        \hline
        Training   & \NumTrainingHealthy{} & \NumTrainingUnhealthy{} & \NumTrainingImages{} \\
        Validation & \textbf{[insert healthy]} & \textbf{[insert unhealthy]} & \NumValidationImages{} \\
        Test       & \textbf{[insert healthy]} & \textbf{[insert unhealthy]} & \NumTestImages{} \\
        \hline
    \end{tabular}
\end{table}

During classifier training, brightness, contrast, and saturation jitter is
applied with probability 0.5 and Gaussian blur is applied with probability
0.3.  Validation and test images receive only deterministic resizing and tensor
conversion.  To prevent the majority class from dominating optimization, the
training loader uses inverse-frequency weighted sampling.  The binary loss is
also parameterized from the training counts using the ratio
$N_{\mathrm{unhealthy}}/N_{\mathrm{healthy}}$.  The classifier is initialized
from ImageNet weights, its final layer is replaced for binary classification,
and it is optimized with Adam for five epochs using a learning rate of
$10^{-4}$ and mini-batches of 32 images.

\subsection{Evaluated Adversarial Attacks}

We study targeted universal perturbations rather than generating a separate
perturbation for every test image.  A single learned trigger is optimized on
the training split, selected using attack success rate (ASR) on the validation
split, and evaluated once on the held-out test split.  The primary threat model
targets unhealthy images and attempts to make the classifier report the
healthy label.  This direction is operationally important because a successful
attack can conceal an abnormal detector condition.  The implementation supports
iterative sign/PGD and momentum-sign updates, DeepFool-style UAP, GD-UAP,
GAP-UAP, HP-UAP, FG-UAP, Robust-UAP, PSP-UAP, and direct Adam optimization.
Perturbations can cover the full image or a localized patch and can be attached
by additive blending or pixel replacement.  For a fair comparison, attacks are
run with the same data partition, target class, perturbation budget, and number
of optimization steps; method-specific hyperparameters are reported with the
corresponding result.

\subsection{Evaluated Defense Methods}

The benchmark includes three input-level defenses.  \emph{Feature
distillation} applies a JPEG-inspired transformation intended to remove
adversarial high-frequency content.  \emph{Diffusion purification} corrupts
and reconstructs the input with a diffusion model trained on the Hydra training
split.  \emph{Feature squeezing} compares model outputs before and after
multiple input reductions and flags inputs whose prediction discrepancy exceeds
a calibrated threshold.  Each defense uses the same pretrained classifier and
held-out test images.  Any calibration required by a defense uses the training
split only, thereby avoiding test-set leakage.

\subsection{Evaluation Protocol}

Clean classification performance is first measured on the unmodified test
set.  For each attack, the trigger is learned from training data, selected on
validation data, and then applied to the eligible test samples.  We report
accuracy, precision, recall, and F1-score for classification.  Attack success
rate is computed over source-class samples that satisfy the attack eligibility
criterion, and the denominator is reported alongside ASR to make comparisons
robust to class imbalance.  Defense evaluation reports performance on both
clean and attacked inputs, together with the mean per-image defense runtime.
This protocol measures both whether conventional computer-vision attacks
transfer to scientific detector images and whether a defense offers sufficient
robustness without compromising the near-real-time monitoring requirement.
